\SourceLeaf{524}Form enthaltenen Reste $=\frac12(p-1)p^{\pi-3}$ ist; für jeden von ihnen gilt jedoch $M=2p$. Die Summe der daraus entstehenden Terme in $\Pi$ wird
\[
\begin{aligned}
&=2p\Sigma r^{\mu pp+kp^3}=2p\Sigma r^{\mu pp}\pdot\Sigma r^{kp^3}\\
&=2p\Sigma r^{\mu pp}\pdot\frac{r^{p^\pi}-1}{r^{p^3}-1}=0
\end{aligned}
\]
vorausgesetzt, dass $\pi>3$. Ebenso wird die dritte, vierte usw. Klasse durch $\mu p^4+kp^5$, $\mu p^6+kp^7$ usw. dargestellt, wobei für $k$ alle ganzen Zahlen von $0$ bis $p^{\pi-5}-1$, $p^{\pi-7}-1$ usw. zu nehmen sind; für diese gilt $M=2pp$, $M=2p^3$ usw. Die Summe der aus der dritten, vierten usw. Klasse entstehenden Terme in $\Pi$ verschwindet, sofern jeweils $\pi>5$, $\pi>7$ usw.. gilt.
\par Daraus folgt für den Fall, dass $\pi$ gerade ist: In $\Pi$ bleiben nur die dem Rest $0$ entsprechenden Terme übrig, die $=1$ sind; für diese gilt aber $M=p^{\frac12\pi}$, sodass die Summe aller Terme in $\Pi$ $=p^{\frac12\pi}$ wird.
\BookRule{.27\textwidth}
